% =============================================================================
% Figura del pipeline experimental — Sección 3.1
%
% Requiere en el preámbulo:
%   \usepackage{tikz}
%   \usetikzlibrary{arrows.meta, positioning, calc}
% =============================================================================
% !TeX root = ../main.tex

\begin{figure}[htbp]
	\centering
	
	\begin{adjustbox}{
			max width=\textwidth,
			max totalheight=\textheight,
			keepaspectratio
		}
		\begin{tikzpicture}[
			node distance=5mm and 6mm,
			font=\small,
			% Bloques generales
			bloque/.style = {
				rectangle,
				rounded corners=2pt,
				draw,
				thick,
				minimum height=8mm,
				align=center,
				inner sep=5pt,
				text width=52mm
			},
			datos/.style = {
				bloque,
				fill=black!5
			},
			ejec/.style = {
				bloque,
				fill=black!12
			},
			analisis/.style = {
				bloque,
				fill=black!7
			},
			salida/.style = {
				bloque,
				fill=white
			},
			% Agentes
			agente/.style = {
				rectangle,
				rounded corners=2pt,
				draw,
				minimum height=9mm,
				align=center,
				inner sep=4pt,
				text width=26mm,
				font=\footnotesize
			},
			ciego/.style = {
				agente,
				very thick,
				fill=black!8
			},
			% Anotaciones
			nota/.style = {
				font=\scriptsize,
				align=left,
				text width=35mm
			},
			% Flechas
			flecha/.style = {
				-{Stealth[length=2mm]},
				thick
			}
			]
			% =====================================================================
			% 1. CONSTRUCCIÓN DEL BANCO DE PRUEBAS
			% =====================================================================
			
			\node[datos] (snap)
			{\textbf{Snapshot financiero}\\
				\footnotesize S\&P 500 · 12/05/2026 · 503 tickers};
			
			\node[datos, below=of snap] (filtro)
			{\textbf{Filtrado y estratificación}\\
				\footnotesize sede en EE.\,UU. · capitalización 5--50\,B\$};
			
			\node[datos, below=of filtro] (arq)
			{\textbf{21 arquetipos}\\
				\footnotesize 7 sectores $\times$ 3 perfiles financieros};
			
			% Ramas: contrafactuales y placebo
			\node[datos, below left=8mm and 3mm of arq] (cestas)
			{\textbf{63 cestas contrafactuales}\\
				\footnotesize control · género · geografía};
			
			\node[datos, below right=8mm and 3mm of arq] (placebo)
			{\textbf{21 parejas placebo}\\
				\footnotesize \textit{placebo\_a} $\leftrightarrow$ \textit{placebo\_b}};
			
			% Anotaciones de tamaño de campaña
			\node[nota, left=4mm of cestas]
			{\textbf{Detección}\\
				6 composiciones\\
				$\times$ 3 niveles\\
				$\times$ 3 semillas\\
				= 54 ejecuciones};
			
			\node[nota, right=4mm of placebo]
			{\textbf{Placebo}\\
				6 composiciones\\
				$\times$ 3 semillas\\
				= 18 ejecuciones};
			
			% =====================================================================
			% 2. EJECUCIÓN EXPERIMENTAL
			% =====================================================================
			
			\node[ejec, below=15mm of arq] (config)
			{\textbf{Configuración experimental}\\
				\footnotesize composición · nivel de instrucción · semilla};
			
			\node[ejec, below=of config] (comite)
			{\textbf{Comité multiagente}\\
				\footnotesize 3 agentes basados en LLM};
			
			% Agentes en paralelo
			\node[agente, below=7mm of comite] (a2)
			{\textbf{Agente 2}\\Visible};
			
			\node[ciego, left=4mm of a2] (a1)
			{\textbf{Agente 1}\\Ciego};
			
			\node[agente, right=4mm of a2] (a3)
			{\textbf{Agente 3}\\Visible};
			
			\node[nota, right=5mm of a3]
			{El agente ciego no recibe\\
				la sede, el CEO ni las\\
				noticias recientes.};
			
			% Punto de convergencia de los tres agentes
			\coordinate (agentes_join) at ($(a2.south)+(0,-5mm)$);
			
			\node[ejec, below=11mm of a2] (genesis)
			{\textbf{Génesis ($t=0$)}\\
				\footnotesize decisión independiente · sin comunicación};
			
			\node[ejec, below=of genesis] (debate)
			{\textbf{Deliberación ($t=1,\ldots,4$)}\\
				\footnotesize intercambio de decisiones y revisión de la asignación};
			
			\node[salida, below=of debate] (registro)
			{\textbf{Registro de resultados}\\
				\footnotesize asignación · recomendación · justificación · metadatos};
			
			% =====================================================================
			% 3. ANÁLISIS
			% =====================================================================
			
			\node[analisis, below=of registro] (deteccion)
			{\textbf{Fase de detección}};
			
			\node[analisis, below=of deteccion] (mitigacion)
			{\textbf{Fase de mitigación}};
			
			% =====================================================================
			% FLECHAS
			% =====================================================================
			
			% Dataset
			\draw[flecha] (snap) -- (filtro);
			\draw[flecha] (filtro) -- (arq);
			\draw[flecha] (arq) -- (cestas);
			\draw[flecha] (arq) -- (placebo);
			
			% Cestas -> configuración
			\draw[flecha] (cestas.south) |- (config.west);
			\draw[flecha] (placebo.south) |- (config.east);
			
			% Ejecución
			\draw[flecha] (config) -- (comite);
			
			% Comité -> agentes en paralelo
			\draw[flecha] (comite.south) -- ++(0,-3mm) -| (a1.north);
			\draw[flecha] (comite.south) -- (a2.north);
			\draw[flecha] (comite.south) -- ++(0,-3mm) -| (a3.north);
			
			% Agentes -> génesis
			\draw[thick] (a1.south) |- (agentes_join);
			\draw[thick] (a2.south) -- (agentes_join);
			\draw[thick] (a3.south) |- (agentes_join);
			\draw[flecha] (agentes_join) -- (genesis.north);
			
			% Génesis -> deliberación -> registro
			\draw[flecha] (genesis) -- (debate);
			\draw[flecha] (debate) -- (registro);
			
			% Análisis
			\draw[flecha] (registro) -- (deteccion);
			\draw[flecha] (deteccion) -- (mitigacion);
			
		\end{tikzpicture}
	\end{adjustbox}
	
	\caption[Pipeline experimental]{Visión general del pipeline experimental,
		desde la construcción del banco de pruebas hasta las dos fases de análisis.
		El primer agente actúa como control interno del diseño: recibe la misma
		información en todas las variantes contrafactuales.}
	\label{fig:pipeline_experimental}
	
\end{figure}